%!TEX TS-program = pdflatex
%!TEX encoding = UTF-8 Unicode
% =============================================================================
%  Cover letter -- Bernoulli submission
% =============================================================================
\documentclass[11pt]{letter}
\usepackage[a4paper,top=2.0cm,bottom=2.0cm,left=2.5cm,right=2.5cm]{geometry}
\usepackage{microtype}
\usepackage[T1]{fontenc}

%\signature{Yuan-chin Ivan Chang}
%\address{Yuan-chin Ivan Chang\\
%  Institute of Statistical Science\\
 % Academia Sinica\\
 % 128 Academia Road, Section~2, Nankang\\
%  Taipei 115, Taiwan\\
%  \texttt{ycchang@stat.sinica.edu.tw}}

\begin{document}

\begin{letter}{The Editors\\ \textit{Bernoulli}\\ Bernoulli Society for
  Mathematical Statistics and Probability}

\opening{Dear Editors,}

I am pleased to submit the enclosed manuscript, \textit{``Martingale
Innovations from Recurrent Networks for Sequential Change-Point
Detection,''} for consideration for publication
in \textit{Bernoulli}.

The contribution is theoretical and deliberately two-tier.  At the first
tier, for operator-norm contractive recurrent networks, the hidden state
admits a Doob decomposition into a predictable component and a
martingale-difference innovation; the transient innovation approaches its
stationary counterpart geometrically, and the accumulated innovations obey
a functional central limit theorem.  A bridge proposition translates this
state-level stabilization into a conditional-drift bound for prediction
residuals under a Lipschitz readout.  At the second tier, for any fitted
predictor whose residuals satisfy explicit conditional-drift and
sub-Gaussian tail conditions, the Innovation CUSUM enjoys a \emph{sharp}
exponential in-control average-run-length lower bound --- at the Lundberg
rate that coincides with the classical i.i.d.\ CUSUM exponent, proved by a
renewal argument rather than a union bound --- together with a logarithmic
detection-delay upper bound.  The two combine into the exponential-alarm /
logarithmic-delay trade-off, which we show is \emph{minimax optimal} in
Lorden's sense, and exactly so under a Gaussian innovation shift; the
guarantee is model-free with respect to the data-generating dynamics though
not assumption-free.  Two further propositions identify when learning is
statistically useful: under a nonlinear conditional mean, fixed linear
residualization carries an irreducible innovation defect that a consistent
learned predictor removes, at a rate that transfers to the monitoring
guarantee.

The empirical studies examine these mechanisms rather than claim a
performance victory: learned innovations are closer to martingale
differences under misspecification, calibration is more stable when the
in-control window is short, and equal-ARL comparisons delimit when the
learned detector is competitive with correctly specified parametric
procedures.  The work sits at the interface of probability and
mathematical statistics, between the residual-based monitoring tradition
(where the model class is specified) and conformal/e-detector approaches
(where exchangeability-type conditions replace it), and is, to our
knowledge, the first to combine recurrent-state innovation theory with
sequential CUSUM calibration for a learned predictor.

In the interest of full transparency, I note that a companion manuscript,
\textit{``Backward Coherence and Hidden-State Stability in Recurrent Neural
Networks: A Quasi-Reverse-Martingale Theory''} (arXiv:2606.08934), develops a
complementary \emph{reverse}-filtration theory for the convergence and stability
of the hidden state itself. The present paper works in the \emph{forward}
filtration, with the innovations; the two manuscripts are self-contained and
neither relies on the other's results. I would be glad to provide the companion
to the editors or referees on request.

The manuscript is original, has not been published previously, and is not under
consideration for publication elsewhere. All proofs, simulation details, and
reproducible code are provided in the accompanying Supplementary Material. I have
no conflicts of interest to declare.

Thank you for your consideration. I would be happy to provide any additional
information the editors may require.

\closing{Sincerely,}
Yuan-chin Ivan Chang\\
ycchang@sinica.edu.tw

\end{letter}
\end{document}
